\section{Polynomial confinement and dense arithmetic profiles}
\label{p14:comp:confinement}

We use the block and alignment conventions of
Section~\ref{p13:conf:section}.  The group argument first treats
primitive triples and nonaligned pairs, then proves the two triple
interfaces stated there.

\begin{lemma}[Polynomial confinement of common divisors]
\label{p13:conf:groups}
Fix \(\tau>0\), \(0<\delta<1/2\), and \(C_0>0\).
Let \(x,y\asymp N\), and let pairwise coprime integers
\(g_a>N^\tau\) satisfy
\[
 g_a\mid x+i_a,\quad g_a\mid y+j_a,\qquad
 |i_a|,|j_a|\le C_0B.
\]
Suppose no line through \((-x,-y)\) contains two
distinct points among the \((i_a,j_a)\).
Then their number, with multiplicity, is
\(O_{\tau,\delta,C_0}(B^{1/2+\delta})=o(B)\).
The same estimate holds after points on exact alignments
have been removed.
\end{lemma}

\begin{proof}
Choose a fixed degree \(d\) such that, for
\(M=(d+1)(d+2)/2\), \((M-1)\tau>d\).
Let \(v_d(s,t)\) list the \(M\) monomials of total degree
at most \(d\), including \(1\).
If \(M\) of the vectors \(v_d(i_a,j_a)\) form a basis,
replace each row of that basis in turn by \(v_d(-x,-y)\).
Subtracting the new row from the other rows shows that
the resulting determinant is divisible by the product
of their \(M-1\) pairwise coprime divisors.
Its absolute value is at most
\[
 C_dN^dB^{d(M-1)}<N^{(M-1)\tau}.
\]
All the determinants vanish.
The vector \(v_d(-x,-y)\) would then belong to the
intersection of the \(M\) coordinate hyperplanes of
this basis, which is \(\{0\}\). Its constant coordinate
is one, a contradiction.

Thus a nonzero polynomial \(F(S,T)\) of degree at most
\(d\) vanishes on all the offset points.
For a geometrically integral component of degree at least
two, Pila's uniform estimate
\cite[Theorem~A]{pila1995} gives
\(O_d(B^{1/2+\delta})\) integer points of height \(O(B)\).
A rationally irreducible component that is not geometrically
integral has only \(O_d(1)\) rational points, by intersecting
two conjugate components.

A line through at least two distinct grid points has a
primitive equation \(aS+bT=c\) with
\(|a|,|b|\ll B\), \(|c|\ll B^2\).
Every group on it divides \(ax+by+c\).
This integer is nonzero by the hypothesis, and has size
\(O(NB)\). Coprimality bounds the number of these groups
by \(O_\tau(1)\). A line through only one grid point
needs no separate point count.
Finally, all groups at one point divide its true value
\(x+i_a=O(N)\); their multiplicity is \(O_\tau(1)\).
There are at most \(d\) components. These observations
prove the bound.
\end{proof}

The groups in this lemma can be composite. Only divisibility
and coprimality are used, and the auxiliary curve is not
enumerated. Its degree is chosen after \(\tau\) and before
\(N\).

\begin{theorem}[Dense primitive triples]
\label{p13:conf:dense-primitive}
For fixed \(\eta,\varepsilon>0\), primitive triples of
true windows with \(\rho_{\mathrm{rel}}(T)\ge\eta B\)
have \(O_{\eta,C,\varepsilon}(N^\varepsilon)\) hosts.
Their homogeneous mass is \(O(N^{-2+\varepsilon})\).
The total absolute centered activity of all primitive
triples is \(O(N^{-3/2+\varepsilon})\).
\end{theorem}

\begin{proof}
Realize the relation space \(R\) linearly on complements
to the local kernels. Its projection into each of the
three blocks is injective, by primitivity.
Apply Lemma~\ref{p1:fixed:recovery} with \(d=2\),
choosing a fixed \(t\) large enough that the loss is at
most \(\eta B/4\). On the resulting space \(R'\),
each nonzero coordinate label has multiplicity at most
\(K=t-1\), and a label has at most \(2t\)
representations by two distinct nonzero labels.

A prime above \(H\) occurs oddly at most once per block.
On the active coordinates of \(R'\), its row has either
two equal nonzero labels or three distinct labels
summing to zero; it cannot have weight one.
Group the primes by their active physical support \(A\),
and let \(G_A\) be their product.
The groups are pairwise coprime and divide all true
values at their support positions.
Every active coordinate belongs to at most
\[
 C_t=K+2tK^2
\]
groups. This grouping uses supports after restriction;
a former ternary row that has become binary is included
as a binary group.

Assign each group to a pair of blocks it meets.
For a primitive triple no such pair has two exact units:
the corresponding even selections would give equal
nonzero square classes in the two blocks.
Their nonzeroness follows from
Lemma~\ref{p1:arith:runge} applied to two distinct values.
Lemma~\ref{p13:conf:groups} therefore shows that only
\(o(B)\) groups exceed \(N^\tau\).
Set every position incident to one of these groups
equal to zero. This loses at most three times their
number of dimensions.
The surviving space has dimension at least \(\eta B/2\)
for large \(N\), and its projections remain injective.

Every still-active position belonged only to original
groups of size at most \(N^\tau\). All its odd primes
above \(H\) were in such groups; hence
\(K_H(n)\le N^{C_t\tau}\).
Each block has at least two active positions, because
its projection contains a nonzero even vector.
Its centre lies in \(\mathcal X_{2C_t\tau}\).
Lemma~\ref{p13:half:small-kernels} bounds the number
of triples by \(N^{6C_t\tau+3\varepsilon'}\).
Choose \(t\) after \(\eta\), and then \(\tau,\varepsilon'\)
after the final margin.

For a primitive triple \(\rho_{\mathrm{rel}}\le L\),
so its homogeneous weight is \(O(N^{-2})\).
For the total activity split at \(\rho_{\mathrm{rel}}=\eta L\).
Below that threshold,
Corollary~\ref{p13:half:connected} gives
\(N^{-3/2+\eta+\varepsilon'}\); above it use the just
proved bound. A primitive coefficient has magnitude
at most its homogeneous weight, whether compatible
or incompatible. Choose \(\eta\) within the final margin.
\end{proof}

\begin{corollary}[Dense nonaligned pairs]
\label{p13:conf:dense-pairs}
For fixed \(\eta,\varepsilon>0\), nonaligned pairs
with \(d_{xy}\ge\eta B\) have \(O(N^\varepsilon)\)
hosts.
\end{corollary}

\begin{proof}
Realize their effective relation space with injective
projections. Regularize with
Lemma~\ref{p1:fixed:recovery}, now with \(d=1\),
losing at most \(\eta B/4\).
Every active prime row has two equal labels.
A coordinate belongs to at most \(K=t-1\) active supports.
Confinement bounds the groups above \(N^\tau\) by \(o(B)\).
Kill both endpoints of every such group.
The remaining positive-dimensional space has two active
positions in each block, with kernels at most \(N^{K\tau}\).
The centres lie in \(\mathcal X_{2K\tau}\).
Their count is at most \(N^{4K\tau+2\varepsilon'}\);
choose \(\tau,\varepsilon'\) sufficiently small.
\end{proof}


\subsection{Proof of Proposition~\ref{p13:conf:near-primitive}}
\label{p14:comp:confinement:near}

\begin{proof}
The kernel of each projection of the relation space
has dimension at most \(\theta B\).
Intersect three complements to these kernels.
The resulting subspace loses at most \(3\theta B\)
dimensions and has all projections injective.

In a given pair, two exact units impose a reduced rational
alignment of height \(O(B)\). Two distinct such alignments
would bound the centres by \(O(B^3)\), which is impossible
at the present scale. The alignment is therefore unique.
If it has \(u\) units, their even-selection space has
dimension \(u-1\); its image loses at most one more
dimension by Lemma~\ref{p1:limits:local-rigidity}.
Thus \(u\le d_{xy}+2\le\theta B+2\).
Kill all the positions of these units in the three pairs,
at a cost of at most \(6\theta B+12\) dimensions.
The remaining space still has macroscopic dimension
and injective projections.

Regularize with a further loss at most \(\eta B/8\).
No pair of surviving offsets lies on an exact alignment,
so the confinement argument in
Theorem~\ref{p13:conf:dense-primitive} applies.
After all positions incident to large groups have been
killed, the centres again belong to
\(\mathcal X_{2C_t\tau}\).
The same count gives \(N^\varepsilon\) hosts.
We never enumerate the chosen subspaces or deletion sets.

Finally,
\(\rho_{\mathrm{rel}}(T)\le d_z+d_{xy}\le L+\theta B\).
The homogeneous weight is at most
\(N^{-2+\theta+o(1)}\), proving the mass estimate.
\end{proof}

\subsection{Proof of Proposition~\ref{p13:conf:strong-profile}}
\label{p14:comp:confinement:strong}

\begin{proof}
Choose the strong pair \(e=\{x,y\}\). Put
\(\delta_z=\dim(S_z\cap(S_x+S_y))
 =\rho_{\mathrm{rel}}(T)-d_{xy}\).
Connectedness gives \(\delta_z>0\), and compatibility gives
\(q(T)=q(e)p_z2^{\delta_z}\).

For \(0<\delta_z<\eta'L\), the conditional count
in Theorem~\ref{p13:half:conditional} gives
\(N^{1/2+\varepsilon'}\) targets of weight at most
\(N^{-1+\eta'+o(1)}\).
Lemma~\ref{p1:arith:dense}, together with
Corollary~\ref{p13:half:pairs} for the independent
part, bounds the total joint mass of strong nonaligned
pairs by \(N^{-1+\varepsilon'}\).
This branch costs \(N^{-3/2+\eta'+O(\varepsilon')}\).

For \(\delta_z\ge\eta'L\), the effective-space form of
Proposition~\ref{p1:fixed:targets}, explained after
\eqref{p13:exact:kappa3}, gives \(N^{\varepsilon'}\)
targets after the fixed pair.
Corollary~\ref{p13:conf:dense-pairs} also gives
\(N^{\varepsilon'}\) source pairs.
Each triple has \(q(T)\ll N^{-3+\sigma+o(1)}\).
Choose \(\eta'\) and the internal margins after
\(\varepsilon\).
\end{proof}
